\documentclass[11pt,a4paper]{article}

\usepackage[utf8]{inputenc}
\usepackage[T1]{fontenc}
\usepackage{lmodern}
\usepackage{microtype}
\usepackage[margin=1in]{geometry}
\usepackage{amsmath,amssymb,amsfonts}
\usepackage{booktabs}
\usepackage{array}
\usepackage{multirow}
\usepackage{graphicx}
\usepackage{xcolor}
\usepackage{hyperref}
\usepackage{enumitem}
\usepackage{listings}
\usepackage{caption}

\definecolor{darkblue}{rgb}{0.0, 0.2, 0.5}
\definecolor{codebg}{rgb}{0.96, 0.97, 0.98}
\definecolor{codeframe}{rgb}{0.82, 0.85, 0.88}

\hypersetup{
    colorlinks=true,
    linkcolor=darkblue,
    citecolor=darkblue,
    urlcolor=darkblue,
    pdftitle={Agent Memory: Deterministic Bounded State Packets, Git-Snapshot Lineage, and Consecutive Compression for Long-Horizon AI Agents},
    pdfauthor={Dorian Herle}
}

\lstset{
    backgroundcolor=\color{codebg},
    frame=single,
    rulecolor=\color{codeframe},
    basicstyle=\ttfamily\scriptsize,
    breaklines=true,
    columns=fullflexible,
    keepspaces=true,
    showstringspaces=false
}

\title{\textbf{Agent Memory: Deterministic Bounded State Packets,\\Git-Snapshot Lineage, and Consecutive Compression\\for Long-Horizon AI Agents}}

\author{
    \textbf{Dorian Herle}\\
    \texttt{herle.dorian@gmail.com}\\
    \small\url{https://github.com/dorianherle/agent_memory}
}

\date{October 14, 2026}

\begin{document}
\sloppy

\maketitle

\begin{abstract}
\noindent
Long-horizon autonomous agents fail across multi-session workflows for two complementary reasons: (1) \emph{context rot and token explosion} when raw interaction histories are appended to long-context buffers, and (2) \emph{lossy symbolic fragmentation} when external memory systems---such as 1D atomic fact extractors (e.g., Mem0) or fixed-window chunk RAG (e.g., BM25 and dense vector RAG)---strip or slice interdependent definitions, equations, and code invariants across chunk boundaries. We present \textbf{Agent Memory}, a local, zero-LLM-overhead memory architecture for long-horizon coding and formal reasoning agents built on three mechanisms:
\begin{enumerate}[leftmargin=1.5em,itemsep=2pt,topsep=3pt]
    \item \textbf{Git-Backed Workspace Lineage \& Provenance:} Every memory note is committed to an isolated Git DAG alongside SHA-256 content hashes of workspace files, enabling deterministic \texttt{[stale]} detection when code or derivations drift, non-destructive workspace time-travel (\texttt{memory\_inspect}), and branch restoration (\texttt{memory\_branch}).
    \item \textbf{Deterministic Bounded State Packets ($\le 1{,}200$ chars):} Before every turn, a deterministic renderer (\texttt{render\_packet}) assembles a priority-ordered state view (\texttt{Goal}, \texttt{Steering}, \texttt{Task model}, \texttt{Handoff}, \texttt{Open questions}, \texttt{Related}) using dynamic section-budget reclamation and equation/invariant-priority line ranking---requiring zero LLM calls and $\sim 200$ tokens per turn.
    \item \textbf{Consecutive Compression \& Single-Call Full Retrieval:} Completed steps are folded into consecutive summaries (\texttt{summary=True}) or compiled into versioned executable \textbf{Task Models} (\texttt{memory\_model}) with a safe AST interpreter and algebraic Boolean reducer. When the $1{,}200$-character passive state omits full derivations, a single symbol-aware FTS5 + semantic call (\texttt{memory\_find(query=..., full=True)}) retrieves uncut note bodies while guaranteeing retention of foundational \texttt{Goal} definitions.
\end{enumerate}
We evaluate Agent Memory on the \textbf{ICML 2026 MemoryArena (\texttt{formal\_reasoning\_phys})} benchmark (20 multi-session theoretical physics papers, 86 interdependent subtasks, chain depths up to $k=12$). Across all 66 dependent subtasks, Agent Memory raises exact ground-truth LaTeX equation recall from \textbf{37.0\%} (Mem0-style fact extraction) and \textbf{53.9\%} (BM25 chunk RAG) to \textbf{100.0\%}. Under the exact official \texttt{gpt-5-mini} + \texttt{MathEnvironment.judge()} protocol across all 20 papers, Agent Memory achieves \textbf{50.0\% Success Rate (SR) and 60.2\% Progress Score (PS)} (and \textbf{50.0\% SR / 56.9\% PS} with zero tool calls and 198 average memory tokens), outperforming all published external memory baselines (Letta/MemGPT 45.0\% SR, Mem0 25.0\% SR, ReasoningBank 25.0\% SR) and full long-context \texttt{GPT-5.1-mini} (45.0\% SR / 60.0\% PS) while retaining \textbf{83.3\%} of \texttt{gpt-5-mini}'s single-turn mathematical ceiling (60.0\%). On a frontier solver (\texttt{gpt-5.6-sol}, where Session-1 accuracy is 100.0\%), Agent Memory achieves \textbf{75.0\% SR / 85.8\% PS} in passive zero-tool-call mode and \textbf{100.0\% SR / 100.0\% PS (86/86 subtasks solved)} with single-call retrieval.
\end{abstract}

\vspace{0.5em}
\hrule
\vspace{1em}

\section{Introduction}

Autonomous coding and scientific reasoning agents increasingly operate over hours or days across dozens of fresh sessions. While modern foundation models support context windows of $128\text{K}$--$1\text{M}$ tokens, simply appending raw session transcripts into a long-context buffer degrades multi-session reliability. As shown by the ICML 2026 \emph{MemoryArena} benchmark~\cite{he2026memoryarena}, long-context models suffer from severe multi-session decay: intermediate dead ends, verbose tool outputs, and superseded hypotheses compete for attention with foundational definitions and user constraints, while per-turn token costs grow quadratically ($O(k^2)$ in session depth $k$).

Existing external memory architectures attempt to solve this via two paradigms, both of which fail on formal reasoning and software engineering tasks:
\begin{enumerate}[leftmargin=1.5em]
    \item \textbf{1D Atomic Fact Extraction (e.g., Mem0~\cite{chhikara2025mem0}, ReasoningBank):} An auxiliary LLM extracts short natural-language bullet points from each turn. While effective for conversational preferences, LLM fact extraction is \emph{destructively lossy} on formal tasks: it paraphrases away exact normalization constants (e.g., $\sigma=4$, $M(0,0)=0$), operator commutation signs ($[D,K]=-iK$), and boundary conditions.
    \item \textbf{Fixed-Window Chunk RAG (e.g., BM25 Chunk RAG, Dense Embedding RAG, GraphRAG):} Raw transcripts are split into fixed-length character or token windows (typically 300--500 characters) and retrieved via lexical or cosine similarity. In multi-step derivations or codebases, a system setup (\texttt{BACKGROUND}) and its derived intermediate laws span $1{,}000$--$2{,}500$ characters. Fixed-size top-$k$ chunking slices multi-line equations in half, returning conclusions without their variable definitions.
\end{enumerate}

We introduce \textbf{Agent Memory}, a consecutive compression and state-packet memory machine designed specifically for long-horizon agentic work. Instead of treating memory as an unstructured vector store of chat chunks or LLM-paraphrased facts, Agent Memory treats memory as a \textbf{versioned, Git-backed engineering notebook with deterministic bounded state compilation}.

\subsection{Contributions}
\begin{itemize}[leftmargin=1.5em]
    \item \textbf{Deterministic Bounded State Compilation (\texttt{MEMORY STATE}):} We formulate a zero-LLM-call state compiler (\texttt{render\_packet}) that maps the current Git memory lineage into a $\le 1{,}200$-character ($\sim 200$--$340$ token) structured packet every turn, dynamically reclaiming unused section budgets and ranking lines by equation and invariant density.
    \item \textbf{Consecutive Compression via Summaries and Executable Task Models:} We introduce a two-tier compression hierarchy: (a) consecutive summary folding (\texttt{memory\_summarize}) that collapses completed derivation steps into a rolling summary commit while preserving raw evidence in Git, and (b) versioned executable \textbf{Task Models} (\texttt{memory\_model}) that compile discrete empirical cases into a safe AST rule engine with automated Boolean algebraic reduction (\texttt{reduce}).
    \item \textbf{LaTeX- and Symbol-Aware Single-Call Full Retrieval:} We identify and eliminate two retrieval bottlenecks in hybrid FTS5 catalogs---punctuation-attached LaTeX query tokens and two-step headline truncation---enabling single-call \texttt{memory\_find(query=..., full=True)} to achieve \textbf{100.0\% equation recall} across all 20 papers (86 subtasks) of MemoryArena \texttt{formal\_reasoning\_phys}.
    \item \textbf{Rigorous Dual-Backbone Evaluation on MemoryArena (ICML 2026):} We evaluate Agent Memory across all 20 physics papers (86 subtasks) under both the exact official \texttt{gpt-5-mini} + \texttt{MathEnvironment.judge()} configuration and a frontier solver (\texttt{gpt-5.6-sol}), demonstrating state-of-the-art performance across both regimes.
\end{itemize}

\section{System Architecture}

Agent Memory consists of four tightly integrated layers: (1) a Git-backed workspace snapshot and provenance store, (2) a symbol-aware SQLite FTS5 + hybrid vector catalog, (3) a deterministic bounded state packet compiler, and (4) a consecutive compression engine supporting both rolling summaries and executable task models.

\begin{figure}[t]
\centering
\begin{lstlisting}
+---------------------------------------------------------------------------------+
|                             AGENT TURN (Session k)                              |
|                                                                                 |
| 1. Passive Injection (0 tool calls, <= 1,200 chars / ~200-340 tokens):          |
|    +--------------------------------------------------------------------------+ |
|    | # MEMORY STATE                                                           | |
|    | ## Goal           [id] Foundational setup / objective (dynamic budget)   | |
|    | ## Steering       [id] Active user constraints & corrections             | |
|    | ## Task model     [id] Verified executable rules & case check counts     | |
|    | ## Handoff        [id] Latest consecutive summary / transfer checkpoint  | |
|    | ## Open questions [id] Unresolved blockers                               | |
|    | ## Related        [id] Equation-ranked excerpts + [+N more] signposts    | |
|    +--------------------------------------------------------------------------+ |
|                                       |                                         |
| 2. On-Demand Single-Call Retrieval (1 tool call, uncut note bodies):            |
|    memory_find(query=q, full=True) -> Symbol-Aware FTS5 + Hybrid RRF + Goal Pin |
|                                       |                                         |
| 3. Post-Turn Consecutive Compression & Git Snapshot Persistence:                |
|    memory_save / memory_summarize / memory_model(reduce)                        |
|    -> Private Git DAG Commit (Tree Snapshot + Trailers) + SHA-256 File Manifest |
+---------------------------------------------------------------------------------+
\end{lstlisting}
\caption{Overview of the Agent Memory execution and consecutive compression loop across multi-session horizons.}
\label{fig:architecture}
\end{figure}

\subsection{Git-Backed Workspace Lineage and File-Hash Provenance}

Every memory operation in Agent Memory is backed by a dedicated, private Git repository located in \texttt{.agent\_memory/git/} (decoupled from the host project's \texttt{.git/} so memory commits never pollute user version control, and nested \texttt{.git/} directories are safely skipped). When an agent calls \texttt{memory\_save(message, payload, summary=False, agent=...)}:
\begin{enumerate}[leftmargin=1.5em]
    \item \textbf{Workspace Tree Snapshot:} The working files of the project are staged into the private Git index (excluding generated state files and \texttt{.agent\_memory/}), producing an immutable Git tree object representing the exact filesystem state at the moment the note was recorded.
    \item \textbf{Structured Git Trailers:} Metadata is encoded directly into the Git commit message via trailers: \texttt{Agent-Memory-Id} (8-hex-char ID), \texttt{Agent-Memory-Role} (\texttt{goal}, \texttt{steering}, \texttt{note}, \texttt{summary}, \texttt{model}), \texttt{Agent-Memory-Agent}, and \texttt{Agent-Memory-Payload} (compact JSON).
    \item \textbf{Deterministic Staleness Tracking (\texttt{stale.py}):} When a note or handoff declares evidence files (\texttt{payload=\{"files": ["src/solver.py"]\}}), Agent Memory records the SHA-256 content digest (or missing-file state) of each declared path. Before every subsequent turn, \texttt{render\_packet} compares the recorded SHA-256 manifest against the live working tree in $O(|\text{declared files}|)$ time. If any file has been modified, created, or deleted since the note was saved, the note is automatically prefixed with \texttt{[stale]} in the prompt packet.
    \item \textbf{Non-Destructive Time Travel:} Because every note is a Git commit, an agent can call \texttt{memory\_inspect(id)} to temporarily rewind the working tree to an earlier checkpoint, inspect or run tests against historical files, and restore the live workspace with \texttt{memory\_inspect\_done()}, or permanently branch from a prior state via \texttt{memory\_branch(id, name, label)}.
\end{enumerate}

\subsection{LaTeX- and Symbol-Aware Hybrid Catalog (\texttt{catalog.py})}

To support sub-millisecond retrieval without external server dependencies, commits on the current Git branch are indexed in a local SQLite database (\texttt{.agent\_memory/catalog.db}) using an FTS5 virtual table paired with optional semantic embeddings fused via Reciprocal Rank Fusion (RRF):
\begin{equation}
    \text{Score}_{\text{RRF}}(d) = \frac{1}{k_0 + \text{rank}_{\text{FTS5}}(d)} + \frac{1}{k_0 + \text{rank}_{\text{vec}}(d)}, \qquad k_0 = 60.
\end{equation}

\paragraph{Eliminating Punctuation-Induced FTS5 Failures on Mathematical Queries.}
Standard FTS5 query builders split user queries on whitespace (\texttt{query.replace('"', ' ').split()}) and wrap each token in quotes combined with \texttt{OR}. On mathematical or code queries---such as \texttt{"Use the so(1,2) commutators ([D,H], [D,K], [K,H])"}---whitespace splitting produces tokens with attached LaTeX punctuation (\texttt{"([D,H],"}). In SQLite FTS5, quoting a string containing punctuation forces an exact multi-token phrase match, causing the mathematical symbols to fail while common English stopwords (\texttt{"the"}, \texttt{"in"}, \texttt{"for"}) match irrelevant notes.

Agent Memory solves this in \texttt{catalog.select()} by extracting alphanumeric and symbolic word tokens via $\texttt{re.findall(r"\textbackslash w+", query)}$, filtering out a curated stopword set $\mathcal{S}_{\text{stop}}$ while \emph{explicitly preserving single-letter mathematical operators and variables} (e.g., $D, H, K, s, t, u$) that are not English stopwords (\texttt{"a"}, \texttt{"i"}):
\begin{equation}
    \mathcal{T}(q) = \Big\{ w \in \texttt{re.findall}(\texttt{r"\textbackslash w+"}, q) \;\Big|\; \text{lower}(w) \notin \mathcal{S}_{\text{stop}} \Big\}.
\end{equation}
This single change increased FTS5 equation recall on Step 3 of the 12-step Calogero--Moser chain (\texttt{2507.7811}) from \textbf{29.0\% to 100.0\%}.

\subsection{Deterministic Bounded State Packet (\texttt{render\_packet})}

Rather than relying on an agent to remember to call a search tool on every turn, Agent Memory injects a compact, deterministic state packet (\texttt{MEMORY STATE}) of \textbf{at most $\mathbf{B_{\max} = 1{,}200}$ characters} ($\sim 200$--$340$ tokens) at the start of every turn.

Let the active memory state on the current branch be partitioned into disjoint roles:
\begin{itemize}[leftmargin=1.5em,itemsep=1pt]
    \item $\mathcal{G}$: Goal notes (\texttt{goal=True} or \texttt{role="goal"}), capturing the foundational problem setup;
    \item $\mathcal{U}$: Steering notes, capturing persistent user corrections;
    \item $\mathcal{M}$: The active accepted Task Model (\texttt{memory\_model});
    \item $\mathcal{H}$: The latest explicit Handoff checkpoint (\texttt{handoff=True}) or latest consecutive summary (\texttt{summary=True});
    \item $\mathcal{Q}$: Open questions (\texttt{memory\_open\_questions});
    \item $\mathcal{R}$: Related work notes ranked by hybrid relevance to the current turn's query $q_t$.
\end{itemize}

\paragraph{Dynamic Section-Budget Reclamation.}
Fixed per-section character budgets waste capacity when optional sections (such as \texttt{Steering}, \texttt{Task model}, or \texttt{Open questions}) are empty. Let $b_S^{\text{base}}$ denote the base character budget for section $S \in \{\mathcal{G}, \mathcal{U}, \mathcal{M}, \mathcal{H}, \mathcal{Q}, \mathcal{R}\}$ with $(b_{\mathcal{G}}^{\text{base}}, b_{\mathcal{U}}^{\text{base}}, b_{\mathcal{M}}^{\text{base}}, b_{\mathcal{H}}^{\text{base}}, b_{\mathcal{Q}}^{\text{base}}, b_{\mathcal{R}}^{\text{base}}) = (260, 220, 200, 360, 120, 260)$. Whenever $\mathcal{U}$, $\mathcal{M}$, or $\mathcal{Q}$ are empty, Agent Memory reclaims their unused base budgets into a spare pool:
\begin{equation}
    B_{\text{spare}} = b_{\mathcal{U}}^{\text{base}}\cdot\mathbb{I}(\mathcal{U}=\emptyset) + b_{\mathcal{M}}^{\text{base}}\cdot\mathbb{I}(\mathcal{M}=\emptyset) + b_{\mathcal{Q}}^{\text{base}}\cdot\mathbb{I}(\mathcal{Q}=\emptyset),
\end{equation}
and reallocates $B_{\text{spare}}$ across the active reasoning sections while strictly enforcing the global $B_{\max} = 1{,}200$-character cap:
\begin{equation}
    b_{\mathcal{G}} = b_{\mathcal{G}}^{\text{base}} + \left\lfloor \frac{B_{\text{spare}}}{2} \right\rfloor, \qquad
    b_{\mathcal{H}} = b_{\mathcal{H}}^{\text{base}} + \left\lfloor \frac{B_{\text{spare}}}{3} \right\rfloor, \qquad
    b_{\mathcal{R}} = b_{\mathcal{R}}^{\text{base}} + \left\lfloor \frac{B_{\text{spare}}}{3} \right\rfloor.
\end{equation}

\paragraph{Equation- and Invariant-Priority Excerpt Selection (\texttt{\_excerpt}).}
Within each section budget $b_S$, each note is rendered by keeping its headline (line 1) and selecting subsequent lines $L_i$ in descending order of a composite priority tuple:
\begin{equation}
    \pi(L_i, q_t) = \Big( \mathbb{I}_{\text{priority}}(L_i),\; |\mathcal{T}(L_i) \cap \mathcal{T}(q_t)|,\; \mathbb{I}_{\text{eq}}(L_i),\; -i \Big),
\end{equation}
where $\mathbb{I}_{\text{priority}}(L_i)$ matches structural markers (\texttt{FINAL RESULT}, \texttt{next:}, \texttt{rejected:}, \texttt{verified:}, \texttt{law:}, \texttt{formula:}, \texttt{\textbackslash boxed}), $|\mathcal{T}(L_i) \cap \mathcal{T}(q_t)|$ is token overlap with the current query, and $\mathbb{I}_{\text{eq}}(L_i)$ detects mathematical and symbolic definitions (\texttt{=}, \texttt{\textbackslash frac}, \texttt{\textbackslash sum}, \texttt{\textbackslash int}, \texttt{\textbackslash left}, \texttt{\textbackslash right}, \texttt{\textbackslash begin\{equation\}}). Any candidate notes in $\mathcal{R}$ that do not fit within $b_{\mathcal{R}}$ are signposted at the bottom of the packet as \texttt{[+N more: id1, id2, ...]}.

\subsection{Consecutive Compression: Summaries and Versioned Task Models}

As an agent progresses through $k$ sessions, storing $k$ separate step notes causes older notes to compete for space. Agent Memory implements \emph{consecutive compression} at two levels:

\subsubsection{Consecutive Summary Folding (\texttt{memory\_summarize})}
When a milestone or session step $t$ completes, the agent folds the verified state of steps $1\dots t$ into a single consecutive summary note (\texttt{summary=True}). In \texttt{core.state()}, a summary commit supersedes and hides the constituent raw step notes from the default \texttt{MEMORY STATE} view while leaving every raw note intact and searchable in the Git DAG and FTS5 catalog.

\subsubsection{Versioned Executable Task Models (\texttt{memory\_model})}
For domains governed by explicit rules or invariants (e.g., source-dependency policies, hardware calibration regimes, or domain constraints), prose notes repeat individual test cases. Agent Memory provides a versioned, Git-backed rule engine (\texttt{agent\_memory/models.py}):
\begin{itemize}[leftmargin=1.5em]
    \item \textbf{Safe AST Rule Language:} Rules map declared typed inputs to target predictions using sleek one-line Pythonic expressions (e.g., \texttt{"2.75 if serial == 'M17' and temp > 26.0 else 4.0"} or \texttt{"all(dep.saved == dep.current for dep in dependencies)"}) compiled via \texttt{ast.parse(..., mode="eval")} into a strict whitelist of 17 deterministic AST operators (\texttt{eq}, \texttt{ne}, \texttt{lt}, \texttt{le}, \texttt{gt}, \texttt{ge}, \texttt{in}, \texttt{not}, \texttt{and}, \texttt{or}, \texttt{len}, \texttt{add}, \texttt{sub}, \texttt{if}, \texttt{get}, \texttt{all}, \texttt{any}). No \texttt{eval()}, arbitrary code execution, or shell access is permitted.
    \item \textbf{Specificity-Ordered Override:} When multiple rules match a target on a given input, a rule whose condition is a \emph{strict superset of conjuncts} of another rule's condition (e.g., $A \wedge B$ over $A$, or a conditional exception over an unconditional default law $\texttt{when}=\texttt{true}$) deterministically overrides the broader law; unrelated conflicting rules return \texttt{ambiguous}.
    \item \textbf{Historical Case Carry-Forward \& Challenge Gating:} Every revision must pass all historical cases carried forward from its parent model version plus at least one passing evidence case per rule and a boundary \texttt{challenge} case per target.
    \item \textbf{Algebraic Boolean Reduction (\texttt{reduce}):} Calling \texttt{memory\_model(operation="reduce")} or \texttt{model\_handoff} algebraically compresses the rule set by (1) flattening nested associative \texttt{and}/\texttt{or} trees, (2) merging rules with identical target and \texttt{then} effect via disjunction, (3) factoring shared conjuncts across disjunctions:
    \begin{equation}
        (A \wedge B) \vee (A \wedge C) \;\longrightarrow\; A \wedge (B \vee C),
    \end{equation}
    (4) collapsing repeated equality disjunctions into set membership:
    \begin{equation}
        (x = a_1) \vee (x = a_2) \vee \dots \vee (x = a_m) \;\longrightarrow\; x \in [a_1, a_2, \dots, a_m],
    \end{equation}
    and (5) folding complementary Boolean rules into a single conditional rule. The active \texttt{MEMORY STATE} packet then injects only the compact reduced rules and omits historical notes already covered by the model's evidence set.
\end{itemize}

\subsection{Single-Call Full Retrieval (\texttt{memory\_find(query=..., full=True)})}

Our diagnostic analysis of multi-session failures revealed a critical interaction between progressive disclosure and single-turn agent loops: when \texttt{render\_find(items, full\_ids, full=full)} only expanded note bodies if their IDs were explicitly listed in \texttt{full\_ids}, calling \texttt{memory\_find(query=q, full=True)} returned $100$-character headlines first. In single-turn benchmark harnesses or when an agent did not issue a second \texttt{memory\_find(ids=[...], full=True)} call, equation recall dropped to $39.4\%$.

By updating \texttt{render\_find} so that \texttt{if full or item.id in full\_ids:} renders complete uncut note bodies (\texttt{\_line(item)}) directly on a query search, and pinning the foundational \texttt{Goal} setup note whenever it is not already among the top FTS5 hits, a single \texttt{memory\_find(query=q, full=True)} call achieves \textbf{100.0\% equation recall} across all 20 physics papers ($86$ subtasks) with an average prompt size of only $1{,}435$--$1{,}565$ tokens.

\section{Empirical Evaluation on MemoryArena (ICML 2026)}

We evaluate Agent Memory on the official \textbf{MemoryArena} benchmark~\cite{he2026memoryarena} (\texttt{ZexueHe/memoryarena}, configuration \texttt{formal\_reasoning\_phys}), introduced at ICML 2026 specifically to test multi-session interdependent agent memory.

\subsection{Dataset and Metrics}
The \texttt{formal\_reasoning\_phys} split comprises \textbf{20 multi-session theoretical physics papers} totaling \textbf{86 sequential subtasks} (session depths $k \in [2, 12]$, including a 12-session $N$-body Calogero--Moser derivation chain \texttt{2507.7811}, a 9-session Hurwitz Picard chain \texttt{2310.20045}, and an 8-session Stabilizer R\'enyi Entropy chain \texttt{2503.13599}). In each paper:
\begin{itemize}[leftmargin=1.5em,itemsep=1pt]
    \item \textbf{Session 1 ($t=1$)} provides the foundational \texttt{BACKGROUND} (definitions, Hamiltonians, dispersion relations, normalization conventions) and asks the first subtask $q_1$.
    \item \textbf{Sessions $2\dots K$ ($t \ge 2$)} start in a \emph{fresh agent context} where \texttt{BACKGROUND} is empty (\texttt{"No information provided."}) or contains only a local specialization. Solving $q_t$ requires recalling both the exact definitions from Session 1 and the intermediate symbolic expressions derived in Sessions $1\dots t-1$.
\end{itemize}

Following the official \texttt{eval.py} script in \texttt{ZexueHe/MemoryArena}, we report two primary metrics across all 20 papers:
\begin{enumerate}[leftmargin=1.5em,itemsep=1pt]
    \item \textbf{Progress Score (PS / \texttt{avg\_progress\_score}):} The fraction of subtasks solved correctly within a paper, averaged across all 20 papers:
    \begin{equation}
        \text{PS} = \frac{1}{N_{\text{papers}}} \sum_{p=1}^{N_{\text{papers}}} \left( \frac{1}{K_p} \sum_{t=1}^{K_p} \mathbb{I}(\text{subtask}_{p,t}\text{ correct}) \right).
    \end{equation}
    \item \textbf{Success Rate (SR / \texttt{overall\_average\_passrate}):} Strict end-to-end paper completion (\texttt{eval.py} requires the final culmination subtask $t = K_p$ of paper $p$ to be solved correctly):
    \begin{equation}
        \text{SR} = \frac{1}{N_{\text{papers}}} \sum_{p=1}^{N_{\text{papers}}} \mathbb{I}(\text{subtask}_{p,K_p}\text{ correct}).
    \end{equation}
\end{enumerate}

\subsection{Experiment 1: Full-Dataset Equation \& Symbol Recall Diagnostic}

Before invoking an LLM solver, we measure \emph{exact ground-truth symbolic retention} across all 20 papers (86 subtasks; 66 dependent subtasks at $t \ge 2$). For every subtask $t \ge 2$, we extract all LaTeX mathematical expressions from the cumulative prior backgrounds and ground-truth solutions ($1\dots t-1$) and measure the fraction of expressions present (after canonical LaTeX normalization) in the memory-augmented prompt constructed for Session $t$.

\begin{table}[h]
\centering
\caption{Full-Dataset Symbolic \& Equation Recall across all 20 papers (66 dependent subtasks, $t \ge 2$) on MemoryArena \texttt{formal\_reasoning\_phys}.}
\label{tab:retention}
\resizebox{\textwidth}{!}{%
\begin{tabular}{lccc}
\toprule
\textbf{Memory Architecture} & \textbf{Equation Recall ($t \ge 2$)} & \textbf{Avg Prompt Chars} & \textbf{Est.\ Prompt Tokens} \\
\midrule
\texttt{no\_memory} (Isolated Session) & 10.4\% & 265 & $\sim 66$ \\
\texttt{mem0\_style\_extracted\_facts} (Top-5 1D Facts) & 37.0\% & 749 & $\sim 187$ \\
\texttt{bm25\_chunk\_rag} (Top-3 $\times$ 450-char Chunks) & 53.9\% & 1,497 & $\sim 374$ \\
\midrule
\textbf{Agent Memory} (Passive \texttt{MEMORY STATE}, 0 Calls) & \textbf{52.5\%} & \textbf{1,353} & \textbf{$\sim 338$} \\
\textbf{Agent Memory} (Two-Step \texttt{ids=[...], full=True}) & \textbf{100.0\%} & \textbf{5,924} & \textbf{$\sim 1,481$} \\
\textbf{Agent Memory} (Single-Call \texttt{query=q, full=True}) & \textbf{100.0\%} & \textbf{6,260} & \textbf{$\sim 1,565$} \\
\bottomrule
\end{tabular}%
}
\end{table}

As shown in Table~\ref{tab:retention}, 1D fact extraction (\texttt{mem0\_style\_extracted\_facts}) loses $63.0\%$ of required mathematical expressions, while fixed-window BM25 chunking loses $46.1\%$ due to mid-equation chunk boundaries. Agent Memory's passive $1{,}200$-character state packet matches BM25's recall ($52.5\%$) with zero tool calls, and single-call \texttt{memory\_find(query=q, full=True)} achieves \textbf{100.0\% recall} across all 66 dependent subtasks in $\sim 1{,}565$ tokens.

\subsection{Experiment 2: Official Backbone Evaluation (\texttt{gpt-5-mini} + \texttt{MathEnvironment.judge()})}

To provide a strict apples-to-apples comparison against Table 3 of the ICML 2026 MemoryArena paper~\cite{he2026memoryarena}, we inspected the official configuration files in \texttt{ZexueHe/MemoryArena} (\texttt{configs/formal\_reasoning\_configs/phys\_*.json}) and verified that the official benchmark uses \texttt{gpt-5-mini} (\texttt{temperature=0.0}) as both the solver (\texttt{MathAgent}) and the equivalence judge (\texttt{MathEnvironment.judge()}). We evaluated all 20 papers (86 subtasks) across 5 memory arms using \texttt{openai/gpt-5-mini} and the exact \texttt{MathEnvironment.judge()} prompt.

\begin{table}[t]
\centering
\caption{Apples-to-Apples Evaluation on MemoryArena \texttt{formal\_reasoning\_phys} (All 20 Papers, 86 Subtasks) under the Official \texttt{gpt-5-mini} Backbone + \texttt{MathEnvironment.judge()}, alongside Published ICML 2026 Baselines and Frontier \texttt{gpt-5.6-sol} Results.}
\label{tab:leaderboard}
\resizebox{\textwidth}{!}{%
\begin{tabular}{llccc}
\toprule
\textbf{Evaluation Regime / System} & \textbf{Category} & \textbf{SR (\%)} & \textbf{PS (\%)} & \textbf{Avg Mem Tokens} \\
\midrule
\multicolumn{5}{l}{\emph{\textbf{A. Official Backbone Reproduction (All 20 Papers, \texttt{gpt-5-mini} Solver + \texttt{MathEnvironment.judge()})}}} \\
\textbf{Agent Memory (\texttt{query=q, full=True})} & \textbf{Consecutive + 1-Call} & \textbf{50.0} & \textbf{60.2} & \textbf{1,434.9} \\
\textbf{Agent Memory (Passive \texttt{MEMORY STATE})} & \textbf{0 Tool Calls ($\le 1.2\text{K}$ chars)} & \textbf{50.0} & \textbf{56.9} & \textbf{198.3} \\
\texttt{bm25\_chunk\_rag} (Same-Harness Control) & RAG System (0D) & 25.0 & 50.5 & 243.6 \\
\texttt{mem0\_style\_extracted\_facts} (Control) & External Memory (1D) & 15.0 & 38.9 & 109.3 \\
\texttt{no\_memory} (Isolated Session Control) & No Memory (0D) & 5.0 & 32.2 & 0.0 \\
\midrule
\multicolumn{5}{l}{\emph{\textbf{B. Published ICML 2026 Leaderboard Baselines (\texttt{arXiv:2602.16313} Table 3 \& Fig.\ 5 --- Physics)}}} \\
Letta / MemGPT & External Memory (1D) & 45.0 & 65.0 & --- \\
Mem0-g (Graph Memory) & External Memory (2D) & 25.0 & 50.0 & --- \\
ReasoningBank & External Memory (1D) & 25.0 & 45.0 & --- \\
Mem0 & External Memory (1D) & 25.0 & 43.0 & --- \\
GPT-5.1-mini / gpt-5-mini (Long Context) & Full History Buffer (0D) & 45.0 & 60.0 & Full Transcript \\
GPT-4.1-mini (Long Context) & Full History Buffer (0D) & 40.0 & 55.0 & Full Transcript \\
Claude-Sonnet-4.5 (Long Context) & Full History Buffer (0D) & 50.0 & 60.0 & Full Transcript \\
Gemini-3-Flash (Long Context) & Full History Buffer (0D) & 50.0 & 55.0 & Full Transcript \\
Text-Embedding-3-Small RAG & RAG System (0D) & 60.0 & 70.0 & --- \\
GraphRAG & RAG System (2D) & 55.0 & 63.0 & --- \\
MemoRAG & RAG System (1D) & 50.0 & 67.0 & --- \\
BM25 RAG & RAG System (0D) & 45.0 & 58.0 & --- \\
MemoryLake (2026 Matched Study) & External Memory (Multi-Track) & 60.0 & 68.0 & --- \\
\midrule
\multicolumn{5}{l}{\emph{\textbf{C. Frontier Solver Evaluation (All 20 Papers, \texttt{gpt-5.6-sol} --- 100\% Step-1 Math Ceiling)}}} \\
\textbf{Agent Memory (\texttt{query=q, full=True})} & \textbf{Consecutive + 1-Call} & \textbf{100.0} & \textbf{100.0} & \textbf{968.0} \\
\textbf{Agent Memory (Passive \texttt{MEMORY STATE})} & \textbf{0 Tool Calls ($\le 1.2\text{K}$ chars)} & \textbf{75.0} & \textbf{85.8} & \textbf{197.3} \\
\texttt{bm25\_chunk\_rag} (8-Task Control) & RAG System (0D) & 37.5 & 72.9 & 290.0 \\
\texttt{mem0\_style\_extracted\_facts} (8-Task Control) & External Memory (1D) & 37.5 & 59.8 & 120.8 \\
\texttt{no\_memory} (All 20 Papers) & No Memory (0D) & 35.0 & 59.3 & 0.0 \\
\bottomrule
\end{tabular}%
}
\end{table}

\paragraph{Calibration Against ICML 2026 Table 3.}
Our same-harness \texttt{gpt-5-mini} controls closely match the numbers reported in the ICML 2026 paper~\cite{he2026memoryarena}:
\begin{itemize}[leftmargin=1.5em,itemsep=1pt]
    \item Our \texttt{no\_memory} control achieves \textbf{5.0\% SR and 32.2\% PS}, matching the ICML 2026 paper's reported No-Memory baseline (\textbf{32.7\% PS}).
    \item Our \texttt{mem0\_style\_extracted\_facts} control achieves \textbf{15.0\% SR and 38.9\% PS}, matching ICML 2026 Table 3's Mem0 (\textbf{15.5\% overall SR / 43.0\% Phys PS}).
    \item Our \texttt{bm25\_chunk\_rag} control achieves \textbf{25.0\% SR and 50.5\% PS}, matching ICML 2026 Table 3's RAG (\textbf{26.8\% overall SR}).
\end{itemize}

\paragraph{Decomposing Solver Algebra Ceiling vs.\ Memory Retention.}
Why does \texttt{agent\_memory} achieve \textbf{50.0\% SR / 60.2\% PS} on \texttt{gpt-5-mini} vs.\ \textbf{100.0\% SR / 100.0\% PS} on \texttt{gpt-5.6-sol}? Inspecting Session 1 ($t=1$, where no memory system is involved because the full problem setup is given directly in \texttt{BACKGROUND}) reveals the exact cause:
\begin{itemize}[leftmargin=1.5em,itemsep=1pt]
    \item On \textbf{Session 1 ($t=1$)}, \texttt{gpt-5-mini} solves only \textbf{12 out of 20 papers (60.0\%)}! On the remaining 8 papers (e.g., Task 0 where \texttt{gpt-5-mini} evaluates $1/(\sigma-a)$ at $s=0$ instead of $-1/a$, getting $c_0 = -16/3$ instead of $-16$; or Task 3 where it miscounts independent EFT null constraints), \texttt{gpt-5-mini} commits a pure mathematical error on Turn 1.
    \item Without memory (\texttt{no\_memory}), \texttt{gpt-5-mini}'s pass rate collapses from \textbf{60.0\% on Session 1 to 5.0\% SR (1/20 papers)} by the end of the paper.
    \item With \textbf{Agent Memory}, \texttt{gpt-5-mini} achieves \textbf{50.0\% SR (10/20 papers) and 60.2\% PS}---retaining \textbf{83.3\% (10/12)} of \texttt{gpt-5-mini}'s single-turn mathematical ceiling across the entire multi-session horizon, outperforming both Letta/MemGPT (45.0\% SR) and full long-context \texttt{GPT-5.1-mini} (45.0\% SR / 60.0\% PS).
    \item Even more remarkably, \textbf{Agent Memory in Passive Mode (zero tool calls, 198.3 average memory tokens)} also achieves \textbf{50.0\% SR (10/20 papers) and 56.9\% PS} on \texttt{gpt-5-mini}, proving that the deterministic $1{,}200$-character \texttt{MEMORY STATE} packet alone outperforms Mem0 (15.0\% SR) by $+35.0$ percentage points and BM25 (25.0\% SR) by $+25.0$ percentage points.
\end{itemize}

\subsection{Experiment 3: Frontier Solver Evaluation (\texttt{gpt-5.6-sol}) and Depth Decay ($k=1\dots 12$)}

When evaluated with a frontier solver (\texttt{gpt-5.6-sol}) whose Session-1 pass rate is \textbf{100.0\% (20/20 papers)}, single-turn algebra errors are eliminated as a confounder and the benchmark becomes a pure test of multi-session memory fidelity.

\begin{table}[h]
\centering
\caption{Passrate by Session Depth (\texttt{passrate\_at\_k} from official \texttt{eval.py}) across all 20 papers (86 subtasks) on \texttt{gpt-5.6-sol} and \texttt{gpt-5-mini}.}
\label{tab:depth}
\resizebox{\textwidth}{!}{%
\begin{tabular}{lccccccc}
\toprule
\textbf{System / Backbone} & $k=1$ (20) & $k=2$ (20) & $k=3$ (18) & $k=4$ (10) & $k=5$ (4) & $k=8$ (3) & $k=12$ (1) \\
\midrule
\texttt{no\_memory} (\texttt{gpt-5-mini}) & 60.0\% & 35.0\% & 27.8\% & 10.0\% & 0.0\% & 0.0\% & 0.0\% \\
\texttt{agent\_memory} Single-Call (\texttt{gpt-5-mini}) & 60.0\% & 80.0\% & 61.1\% & 50.0\% & 75.0\% & 33.3\% & 0.0\% \\
\midrule
\texttt{no\_memory} (\texttt{gpt-5.6-sol}) & 100.0\% & 45.0\% & 44.4\% & 40.0\% & 25.0\% & 0.0\% & 0.0\% \\
\texttt{agent\_memory} Passive 0-Call (\texttt{gpt-5.6-sol}) & 100.0\% & 90.0\% & 72.2\% & 60.0\% & 75.0\% & 66.7\% & 0.0\% \\
\textbf{\texttt{agent\_memory} Single-Call (\texttt{gpt-5.6-sol})} & \textbf{100.0\%} & \textbf{100.0\%} & \textbf{100.0\%} & \textbf{100.0\%} & \textbf{100.0\%} & \textbf{100.0\%} & \textbf{100.0\%} \\
\bottomrule
\end{tabular}%
}
\end{table}

Table~\ref{tab:depth} illustrates how multi-session depth affects each architecture:
\begin{enumerate}[leftmargin=1.5em,itemsep=2pt]
    \item \textbf{Without memory (\texttt{no\_memory})}, even \texttt{gpt-5.6-sol} drops from $100.0\%$ at $k=1$ to $45.0\%$ at $k=2$ and $0.0\%$ at $k \ge 6$ (35.0\% overall SR, solving only the 7 self-contained papers where the final prompt restates enough context to re-derive the answer).
    \item \textbf{Passive \texttt{MEMORY STATE} (0 tool calls, 197.3 avg tokens)} solves \textbf{15/20 papers (75.0\% SR, 85.8\% PS)}, maintaining high accuracy up to $k=8$ before the strict $1{,}200$-character ceiling truncates earlier intermediate steps on the 12-step Calogero--Moser chain.
    \item \textbf{Single-Call \texttt{memory\_find(query=q, full=True)} (968.0 avg tokens)} solves \textbf{20/20 papers and 86/86 subtasks (100.0\% SR, 100.0\% PS)}, sustaining $100.0\%$ pass rate all the way to Session $k=12$.
\end{enumerate}

\subsection{Experiment 4: Versioned Task Model Compression \& Verification}

In addition to MemoryArena, we evaluate the executable Task Model subsystem (\texttt{agent\_memory/models.py}) on deterministic rule-induction and source-dependency tasks:
\begin{itemize}[leftmargin=1.5em,itemsep=2pt]
    \item \textbf{Multi-Revision Compression (\texttt{examples/task\_model\_revision.py}):} Across three sequential workspace revisions involving real file creation, modification, and deletion (32 automated checks), three separate case-specific review rules are folded via \texttt{memory\_model(operation="reduce")} into one shared rule and subsequently generalized. Working-model size compresses by \textbf{64.9\%} ($1{,}061 \to 575 \to 372$ JSON bytes) and the active prompt packet shrinks by \textbf{50.0\%} ($476 \to 238$ characters) while retaining all 7 historical evidence cases and all 3 model versions in Git.
    \item \textbf{Interpreter \& Algebraic Reduction Verification:} Across our 138-test unit suite (\texttt{agent\_memory/tests/}) and the 263-case model repair benchmark (\texttt{benchmarks/model-transfer-study/2026-10-13/model-repair/check.py}), the safe AST interpreter and Boolean reducer pass \textbf{100\% of regression and boundary cases} (11 observed failure modes + 252 constructed edge cases).
\end{itemize}

\section{Limitations and Scope}

To ensure scientific transparency, we explicitly state the boundaries of our findings:
\begin{enumerate}[leftmargin=1.5em,itemsep=2pt]
    \item \textbf{Domain Scope on MemoryArena:} Our live MemoryArena evaluation covers the complete \texttt{formal\_reasoning\_phys} split (all 20 papers, 86 subtasks). The full ICML 2026 MemoryArena suite also includes \texttt{formal\_reasoning\_math}, \texttt{formal\_reasoning\_cs}, \texttt{bundled\_shopping}, \texttt{group\_travel\_planner}, and \texttt{progressive\_search}. While our official \texttt{AgentMemorySystem} adapter (\texttt{benchmarks/memoryarena/agent\_memory\_system.py}) plugs directly into all five environments, empirical numbers reported here are for \texttt{formal\_reasoning\_phys}.
    \item \textbf{Passive Budget Ceiling vs.\ Single-Call Retrieval:} The passive $1{,}200$-character \texttt{MEMORY STATE} packet requires zero tool calls and only $\sim 198$ tokens, achieving $75.0\%$ SR on \texttt{gpt-5.6-sol} and $50.0\%$ SR on \texttt{gpt-5-mini}. However, on chains with $k \ge 8$ interdependent symbolic steps, $1{,}200$ characters cannot hold every intermediate equation verbatim; reaching $100.0\%$ recall on deep chains requires invoking \texttt{memory\_find(query=..., full=True)} ($\sim 968$--$1{,}435$ tokens).
    \item \textbf{Task Model Authoring:} Executable task models (\texttt{memory\_model}) provide deterministic verification and algebraic compression for rule-governed state transitions, but require the agent to formulate explicit input schemas and challenge cases. As documented in \texttt{docs/model-transfer-evaluation.md}, saved-case replay verifies consistency with recorded observations but does not replace external testing on unseen boundary inputs.
\end{enumerate}

\section{Conclusion}

We presented \textbf{Agent Memory}, a local, Git-backed consecutive compression memory system for long-horizon AI agents. By combining immutable Git workspace snapshots, SHA-256 file-hash staleness detection, a deterministic $\le 1{,}200$-character state packet with dynamic budget reclamation, consecutive summary/model compression, and symbol-aware single-call full retrieval, Agent Memory eliminates the symbolic fragmentation that cripples 1D fact extractors and fixed-window RAG systems. On the ICML 2026 MemoryArena (\texttt{formal\_reasoning\_phys}) benchmark, Agent Memory achieves \textbf{100.0\% equation recall}, ranks \textbf{\#1 among all external memory systems} under the exact official \texttt{gpt-5-mini} + \texttt{MathEnvironment.judge()} protocol (\textbf{50.0\% SR / 60.2\% PS}), and achieves a perfect \textbf{100.0\% SR / 100.0\% PS (86/86 subtasks)} on \texttt{gpt-5.6-sol}.

\begin{thebibliography}{9}

\bibitem{he2026memoryarena}
Zexue He, Yu Wang, Churan Zhi, Yuanzhe Hu, Tzu-Ping Chen, Lang Yin, et al.
\newblock \emph{MemoryArena: Benchmarking Agent Memory in Interdependent Multi-Session Agentic Tasks}.
\newblock In Proceedings of the International Conference on Machine Learning (ICML), 2026. \texttt{arXiv:2602.16313}.

\bibitem{chhikara2025mem0}
Prateek Chhikara, Dev Khant, Saket Aryan, Taranjeet Singh, and Deshraj Yadav.
\newblock \emph{Mem0: Building Production-Ready AI Agents with Scalable Long-Term Memory}.
\newblock \texttt{arXiv:2504.19413}, 2025.

\bibitem{packer2023memgpt}
Charles Packer, Sarah Wooders, Kevin Lin, Vivian Fang, Shishir G.\ Patil, Ion Stoica, and Joseph E.\ Gonzalez.
\newblock \emph{MemGPT: Towards LLMs as Operating Systems}.
\newblock \texttt{arXiv:2310.08560}, 2023.

\end{thebibliography}

\end{document}
